\chapter{Introduzione}
\label{sec:introduction}

% Source notes: Project/docs/implementation_plan.md §1–2 (research questions), checkpoint_2026-10-10.md (results).

\section{Motivazione}
\label{sec:intro-motivation}

La generazione aumentata dal recupero (RAG) è diventata il modo standard per consentire a un modello linguistico di
rispondere a domande su documenti che non ha mai visto durante l'addestramento: un sistema di recupero seleziona i
passaggi pertinenti e il modello scrive una risposta tenendo tali passaggi nel proprio contesto~\cite{lewis2020rag}.
Poiché modelli aperti compatti e runtime efficienti girano ormai su hardware di consumo~\cite{llamacpp,mlx2023}, l'intera
pipeline può restare sul dispositivo dell'utente. Ciò è rilevante ogni volta che i documenti sono privati, che
l'assistente deve funzionare senza rete o che il suo comportamento non deve cambiare a seguito di un aggiornamento del
fornitore.

Il recupero non elimina l'allucinazione. Le risposte contengono ancora frasi che i passaggi recuperati non supportano, o
che li contraddicono~\cite{niu2024ragtruth,huang2025hallucination}, e non si può pretendere che un utente controlli ogni
frase rispetto alla sua fonte. Un assistente locale affidabile richiede quindi un verificatore che controlli ogni
risposta, frase per frase, sulla stessa macchina che l'ha generata. Un simile verificatore compete con il generatore per
la stessa memoria, lo stesso tempo e la stessa batteria. La soluzione più diffusa, cioè chiedere a un modello linguistico
di giudicare la risposta, è costosa proprio là dove un laptop è più debole: un giudice scrive i propri verdetti token per
token, e su hardware locale ogni token generato costa molto più di un token letto. I giudici di piccole dimensioni,
inoltre, violano il formato di output richiesto, e la confidenza che esprimono è raramente utilizzabile come probabilità.

\section{Modelli System One}
\label{sec:intro-systemone}

Un modo diverso di usare un modello linguistico per prendere decisioni ha attirato di recente l'attenzione con il nome di
\emph{modelli System One}~\cite{typesafe2026jev}. Anziché generare una risposta, un modello di questo tipo legge una sola
volta la domanda e il contesto e restituisce una probabilità per ciascuna risposta ammissibile, ad esempio leggendo i
logit delle lettere delle opzioni alla fine del prompt~\cite{semif,rizzoflow}. Non viene generato alcun testo, l'output è
sempre valido e molte domande sullo stesso contesto possono condividere un unico passaggio su di esso. L'idea si adatta
bene alla verifica: una risposta di $K$ frasi pone $K$ domande chiuse sullo stesso documento, e una probabilità per
ciascuna frase può essere sottoposta a soglia e calibrata.

Le evidenze a favore dei modelli System One, tuttavia, provengono per lo più dai loro stessi autori, e le
implementazioni aperte sono sistemi generalisti che rispondono a qualunque domanda tipizzata. Restano aperte diverse
questioni: se l'approccio sia vantaggioso per un compito specifico, come si collochi un modello specializzato in quel
compito rispetto a uno generalista, ai giudici generativi e ai verificatori dedicati, e da dove derivi effettivamente
l'eventuale vantaggio.

\section{Domande di ricerca}
\label{sec:intro-rq}

Questa tesi progetta, addestra e valuta GroundOne, un modello System One compatto specializzato in un'unica decisione:
stabilire, per ogni frase di una risposta RAG, se il documento recuperato la supporta, la contraddice o non la menziona.
Il lavoro è guidato da cinque domande di ricerca.

\begin{description}
\item[RQ1, specializzazione.] Un modello System One sottoposto a fine-tuning, con 0.5B o 3B parametri, può eguagliare
modelli zero-shot più grandi e gli encoder specializzati nella verifica della groundedness (aderenza al documento) a
livello di frase?
\item[RQ2, lettura dei logit (readout) contro generazione.] Quale contributo apporta il meccanismo System One a parità di dimensione del
modello, rispetto ai giudici generativi che scrivono i propri verdetti e rispetto allo stesso modello che decide in base
al token che genererebbe, e come cresce la differenza con il numero di decisioni per risposta?
\item[RQ3, calibrazione.] Le probabilità sono utilizzabili come confidenza, per fissare soglie e per rinviare le
decisioni incerte, e come si confrontano con le probabilità che un giudice generativo è in grado di fornire?
\item[RQ4, costo sul dispositivo.] Il verificatore rientra nel budget di memoria, tempo ed energia di un laptop da
16\,GB, e come si confronta il suo costo con quello di un giudice sulla stessa macchina?
\item[RQ5, oltre un punteggio binario.] Che cosa offre una decisione a tre classi per ogni frase che un verificatore
binario non offre?
\end{description}

\section{Contributi}
\label{sec:intro-contributions}

La tesi apporta quattro contributi.
\begin{enumerate}
\item \textbf{GroundOne}, un modello System One specialista e aperto per la verifica della groundedness a livello di
frase, eseguito interamente su un laptop. Legge tutti i verdetti di una risposta in un solo passaggio in avanti, grazie a
un marcatore posto dopo ogni frase, ed è distribuito come pacchetto installabile con un'interfaccia a riga di comando e
un'API HTTP locale, che include un'interfaccia compatibile con i client System One esistenti.
\item \textbf{Un dataset orientato al compito}, costruito a partire da annotazioni umane pubbliche, con tre insiemi di
test validati senza nuove annotazioni e congelati prima dell'addestramento di qualsiasi modello, e con il leakage
controllato a livello di documento.
\item \textbf{Una scomposizione controllata dei guadagni.} Il confronto separa il contributo della specializzazione,
del readout ristretto e della dimensione del modello, rispetto a indizi che non verificano mai una
frase, modelli non addestrati, giudici generativi con e senza probabilità, un modello System One generalista ed encoder
specializzati, tutti sullo stesso hardware e con intervalli di cluster bootstrap.
\item \textbf{Uno studio di calibrazione, tempo ed energia}: come le probabilità possano essere calibrate e tradotte in
etichette coerenti, quanto aiutino a rinviare le decisioni incerte, e come il tempo e l'energia per verifica crescano
con la lunghezza della risposta.
\end{enumerate}

In sintesi, i risultati mostrano che la specializzazione funziona nel dominio RAG: sull'insieme di test RAGTruth
GroundOne raggiunge un'accuratezza bilanciata a tre classi di 0.769 con 0.5B parametri e di 0.832 con 3B, superiore a
quella degli encoder specializzati, nettamente superiore a quella di un modello System One generalista otto volte più
grande e superiore a quella di un modello decisionale generalista della stessa dimensione; su LLM-AggreFact, inoltre, il
modello da 3B è alla pari con il miglior encoder di piccole dimensioni.
Verifica le risposte lunghe circa due volte più velocemente di un giudice sintetico e sei volte più velocemente di un
giudice verboso della stessa dimensione, con un output sempre valido, probabilità molto migliori e un minore consumo di
energia. La restrizione del readout alle risposte ammissibili aumenta da sola l'accuratezza. I limiti sono altrettanto
chiari: il trasferimento a un dominio diverso è debole, e lì il modello generalista risulta migliore; le allucinazioni
nascoste in frasi che per il resto ricopiano la fonte restano difficili da individuare; la ricetta di addestramento,
infine, non è del tutto robusta rispetto al seed casuale.

\section{Struttura della tesi}
\label{sec:intro-structure}

Il Capitolo~\ref{sec:background} passa in rassegna l'inferenza dei transformer, l'inferenza locale, i modelli System One,
l'allucinazione nella generazione aumentata dal recupero, i metodi di verifica e la calibrazione. Il
Capitolo~\ref{sec:method} descrive GroundOne: la decisione, il readout in un solo passaggio, l'obiettivo di
addestramento, la calibrazione e l'output servito. Il Capitolo~\ref{ch:data} definisce le etichette e descrive come gli
insiemi di addestramento e di test sono stati costruiti, verificati e congelati. Il Capitolo~\ref{sec:setup} presenta i
modelli addestrati, le baseline, le metriche e il protocollo di misura. Il Capitolo~\ref{sec:results} riporta i
risultati, il Capitolo~\ref{sec:discussion} ne discute il significato e i limiti, e il Capitolo~\ref{sec:conclusion}
conclude. L'Appendice~\ref{app:examples} mostra l'output di ciascun sistema su una stessa risposta.
